\documentclass[12pt,a4paper]{article}

\usepackage[margin=1in]{geometry}
\usepackage{setspace}
\usepackage{titlesec}
\usepackage{enumitem}
\usepackage{hyperref}

\hypersetup{colorlinks=true,linkcolor=blue,urlcolor=blue}
\setstretch{1.15}

\titleformat{\section}{\large\bfseries}{\thesection.}{0.5em}{}

\begin{document}

\begin{center}
    {\LARGE \textbf{CampusCycle}}\\[0.3cm]
    {\Large \textbf{Smart Campus Bicycle Sharing and Mobility Management System}}\\[0.6cm]
    \textbf{Project Proposal}\\[0.7cm]
    \textbf{Team Members:}\\[0.2cm]
    Guneet Singh Nagi\\
    Prabhnoor Oberoi\\[0.5cm]
    \textbf{Department:} Computer Science and Engineering
\end{center}

\section{Introduction}

CampusCycle is a smart bicycle-sharing platform designed specifically for university campuses. The system allows students to locate, reserve, unlock, use, and return shared bicycles through a web application. Administrators can manage bicycles, stations, rentals, maintenance, and usage analytics.

The project extends beyond basic rental functionality by incorporating QR-based bicycle identification, geofencing, maintenance reporting, usage analytics, and demand prediction. These features aim to make short-distance transportation within a large campus more convenient, affordable, sustainable, and manageable.

\section{Problem Statement}

Large university campuses often require students to travel considerable distances between hostels, academic blocks, libraries, cafeterias, sports facilities, and administrative buildings. Although bicycles can solve this last-mile mobility problem, a shared bicycle system requires efficient tracking and management.

A conventional rental system does not adequately address problems such as uneven bicycle distribution, unauthorized use outside the campus, maintenance issues, and difficulty in predicting where bicycles will be needed.

Therefore, the proposed system aims to develop a centralized smart campus bicycle-sharing platform that manages the bicycle fleet while using location and historical usage data to improve availability and operational efficiency.

\section{Objectives}

The main objectives of CampusCycle are:

\begin{itemize}[leftmargin=*]
    \item Provide students with a simple platform to find and rent campus bicycles.
    \item Implement QR-based identification for individual bicycles.
    \item Display bicycle and docking-station availability on an interactive campus map.
    \item Maintain complete rental and return records.
    \item Use geofencing to detect bicycles leaving the permitted campus area.
    \item Allow students to report bicycle faults and maintenance problems.
    \item Provide administrators with fleet, usage, and maintenance analytics.
    \item Analyze historical rental data to predict bicycle demand at different stations.
    \item Recommend redistribution of bicycles between stations based on predicted demand.
\end{itemize}

\section{Proposed System}

The system will contain two major interfaces: a student application and an administrator dashboard.

\subsection{Student Application}

Students will be able to:

\begin{itemize}[leftmargin=*]
    \item Register and log in using their university account.
    \item View nearby bicycle stations and available bicycles.
    \item Scan a QR code to identify a bicycle.
    \item Reserve and start a bicycle rental.
    \item View the active ride and rental duration.
    \item Return the bicycle at an approved station.
    \item View previous rides and distance travelled.
    \item Report mechanical or safety issues.
    \item View personal mobility statistics and sustainability contributions.
\end{itemize}

\subsection{Administrator Dashboard}

Administrators will be able to:

\begin{itemize}[leftmargin=*]
    \item Add, remove, and update bicycles and stations.
    \item Monitor active rentals.
    \item Track bicycle availability and locations.
    \item View maintenance reports and bicycle status.
    \item Detect geofence violations.
    \item Analyze station-wise and time-wise usage.
    \item View predicted demand for upcoming periods.
    \item Receive recommendations for bicycle redistribution.
\end{itemize}

\section{Key Features}

\subsection{QR-Based Bicycle Identification}

Each bicycle will have a unique QR code. Scanning the code allows the backend to identify the bicycle and verify whether it is available, reserved, currently rented, or under maintenance.

\subsection{Campus Geofencing}

A virtual boundary will be defined around the campus. In an extended implementation with GPS-enabled hardware, the system can detect when a bicycle crosses this boundary and generate an alert for the administrator. For the software prototype, location data can be simulated or obtained from a mobile device.

\subsection{Interactive Campus Map}

The application will display stations and their current availability. Stations can be represented using status indicators such as available, low availability, and unavailable.

\subsection{Maintenance Management}

Students can report issues such as damaged brakes, flat tyres, or faulty lights. Repeated reports can increase the priority of a maintenance request, and bicycles marked unsafe can automatically become unavailable for new rentals.

\subsection{Demand Prediction}

Historical rental information will be analyzed using machine-learning techniques. Features may include station, time of day, day of week, previous demand, and number of active rentals. A regression or tree-based model will predict future demand.

\subsection{Smart Redistribution}

Predicted demand will be compared with current station capacity. The system can recommend moving bicycles from stations with excess supply to stations where demand is expected to be high.

\section{System Architecture}

The proposed architecture consists of:

\begin{enumerate}[leftmargin=*]
    \item \textbf{Frontend:} React.js with Tailwind CSS for the student and administrator interfaces.
    \item \textbf{Backend:} Node.js and Express.js for authentication, rentals, bicycle management, maintenance, and REST APIs.
    \item \textbf{Database:} MySQL for users, bicycles, stations, rentals, locations, maintenance records, and analytics data.
    \item \textbf{Machine Learning Module:} Python with Scikit-learn for demand prediction and redistribution recommendations.
    \item \textbf{Mapping:} Leaflet/OpenStreetMap or another suitable mapping service for the campus map.
    \item \textbf{Optional Hardware Extension:} ESP32 with GPS and a lock mechanism for demonstrating real-time bicycle tracking.
\end{enumerate}

\section{Technology Stack}

\begin{center}
\begin{tabular}{ll}
\textbf{Component} & \textbf{Technology} \\
\hline
Frontend & React.js, Tailwind CSS \\
Backend & Node.js, Express.js \\
Database & MySQL \\
Machine Learning & Python, Scikit-learn, Pandas \\
Maps & Leaflet / OpenStreetMap \\
Authentication & JWT / University Authentication \\
Version Control & Git and GitHub \\
Optional Hardware & ESP32, GPS module, electronic lock \\
\end{tabular}
\end{center}

\section{Expected Outcome}

At the end of the project, a working prototype of CampusCycle will allow a student to discover an available bicycle, scan its QR code, start a rental, view the ride status, and return the bicycle at an approved station. Administrators will have a dashboard for fleet and maintenance management.

The system will additionally demonstrate analytics and a machine-learning model capable of predicting station-level bicycle demand. The prototype will provide a foundation that could later be extended to a real campus deployment with GPS trackers, electronic locks, payment integration, and mobile applications.

\section{Project Scope}

The minimum viable product will include authentication, bicycle and station management, QR-based identification, rental/return workflows, campus map visualization, maintenance reporting, and an administrator dashboard.

Advanced features such as real-time GPS tracking, electronic locks, demand prediction, and smart redistribution will be implemented as higher-level modules depending on available development time and hardware availability.

\section{Conclusion}

CampusCycle combines web development, databases, geospatial concepts, machine learning, and optional IoT hardware into a practical campus mobility solution. The project addresses a real transportation problem while remaining feasible as a two-person semester project. Its modular design also allows the basic bicycle-sharing system to be developed first and advanced intelligence and hardware features to be added incrementally.

\end{document}
